\documentclass[11pt]{article}
\usepackage[margin=1in]{geometry}
\usepackage{amsmath,amssymb,amsthm,mathtools}
\usepackage{graphicx}
\usepackage{microtype}
\usepackage[hidelinks]{hyperref}
\usepackage{xcolor}

\newtheorem{theorem}{Theorem}
\newtheorem{lemma}{Lemma}
\newtheorem{remark}{Remark}
\newcommand{\K}{\mathbb K}
\newcommand{\Lcal}{\mathcal L}
\newcommand{\Piabs}{\Pi}
\newcommand{\ellw}[2]{\ell_{#1}^{w}(#2)}
\newcommand{\norm}[1]{\left\lVert #1\right\rVert}
\newcommand{\ip}[2]{\left\langle #1,#2\right\rangle}

\title{Complete resolution of equation (4):\\a sharp jump at $q=1$}
\author{Autonomous proof candidate for specialist verification}
\date{9 August 2026}

\begin{document}
\maketitle

\begin{abstract}
Let $1<p<2$ and $1\le q\le2$.  We resolve equation~(4) in
Blasco--Botelho--Pellegrino--Rueda: every continuous bilinear form on
$\ell_p\times\ell_q$ is absolutely $(p;2,1)$-summing if and only if $q=1$.
At $q=1$ we prove the stronger $(1;2,1)$ conclusion using Grothendieck
factorization, the little Grothendieck theorem, and trace ideals.  For every
$1<q\le2$ we give an explicit compact-operator counterexample built from
uniformly complemented Rademacher Euclidean blocks.  The literature search
supporting novelty was bounded, so the result should be treated as a full
candidate pending specialist review.
\end{abstract}

\section{Source question and answer}

For a Banach space $E$ and a finite family $(z_j)_{j=1}^m\subset E$, write
\[
 \norm{(z_j)_{j=1}^m}_{\ell_t^w(E)}
 =\sup_{\varphi\in B_{E^*}}
   \left(\sum_{j=1}^m |\varphi(z_j)|^t\right)^{1/t}.
\]
A continuous bilinear form $A:E\times F\to\K$ is absolutely
$(u;r,s)$-summing if a constant $C$ exists such that
\[
 \left(\sum_{j=1}^m|A(x_j,y_j)|^u\right)^{1/u}
 \le C\norm{(x_j)}_{\ell_r^w(E)}\norm{(y_j)}_{\ell_s^w(F)}                 \tag{1}
\]
for all finite families.  The least such constant is denoted
$\pi_{(u;r,s)}(A)$.

On page 9 of the source paper~\cite{BBPR}, after proving nearby coincidence
results, the authors ask whether, for $1<p<2$ and $1\le q\le2$,
\[
 \Lcal(\ell_p,\ell_q;\K)
 =\Piabs_{(p;2,1)}(\ell_p,\ell_q;\K).                                  \tag{2}
\]
The source text is reproduced below.

\begin{center}
 \includegraphics[width=.91\textwidth]{figures/open_problem_crop.png}
\end{center}

Our answer is a sharp endpoint classification.

\begin{theorem}[Full resolution of (2)]\label{thm:main}
Let $1<p<2$ and $1\le q\le2$.  Then
\[
 \Lcal(\ell_p,\ell_q;\K)
 =\Piabs_{(p;2,1)}(\ell_p,\ell_q;\K)
 \quad\Longleftrightarrow\quad q=1.
\]
More precisely:
\begin{enumerate}
\item every $A\in\Lcal(\ell_p,\ell_1;\K)$ is absolutely $(1;2,1)$-summing;
\item for every $1<q\le2$, there is a compact operator
      $T:\ell_q\to\ell_{p'}$ such that
      $A_T(x,y)=\ip{Ty}{x}$ is not absolutely $(p;2,1)$-summing.
\end{enumerate}
\end{theorem}

\section{Proof intuition}

At $q=1$, a weakly $\ell_1$ family is encoded by an operator
$R:\ell_\infty^m\to\ell_1$.  Grothendieck's factorization theorem factors
$R$ through a Hilbert space.  One resulting column operator is
Hilbert--Schmidt by averaging over signs.  The little Grothendieck theorem
makes the other constructed operator Hilbert--Schmidt.  Their product is
trace class, and its diagonal is exactly $(A(x_j,y_j))_j$.  This gives an
$\ell_1$ estimate, stronger than the required $\ell_p$ estimate.

When $q>1$, both $\ell_q$ and $\ell_{p'}$ contain uniformly complemented
Euclidean spaces of every finite dimension.  We use explicit Rademacher
embeddings and projections.  A block-diagonal operator which is a slowly
decaying scalar multiple of the identity on these Euclidean blocks is
compact, yet carefully chosen diagonal test sequences force the quotient in
(1) to diverge.

\section{The positive endpoint \texorpdfstring{$q=1$}{q=1}}

We use two standard forms of Grothendieck's theorem; see, for example,
Pisier's survey~\cite{Pisier}.  First, every operator
$R:\ell_\infty^m\to\ell_1$ admits a Hilbert-space factorization
\[
 \ell_\infty^m\xrightarrow{J}H\xrightarrow{V}\ell_1,
 \qquad R=VJ,\qquad \norm V\norm J\le K_G\norm R.                       \tag{3}
\]
Second, the little Grothendieck theorem says that every
$S:\ell_\infty\to H$ is $2$-summing and
\[
 \pi_2(S)\le K_G\norm S.                                               \tag{4}
\]
The constants may depend on the scalar field and their exact values are
irrelevant.

Fix $A\in\Lcal(\ell_p,\ell_1;\K)$, and let
$T:\ell_1\to\ell_{p'}$ be its associated operator:
$A(x,y)=\ip{Ty}{x}$.  Given finite families $(x_j)_{j=1}^m\subset\ell_p$ and
$(y_j)_{j=1}^m\subset\ell_1$, define
\[
 X:\ell_2^m\to\ell_p,\quad Xe_j=x_j,
 \qquad
 R:\ell_\infty^m\to\ell_1,\quad Ra=\sum_{j=1}^m a_jy_j.
\]
Then
\[
 \norm X=\norm{(x_j)}_{\ell_2^w(\ell_p)},\qquad
 \norm R=\norm{(y_j)}_{\ell_1^w(\ell_1)}.                              \tag{5}
\]
Choose the factorization (3), and introduce
\[
 U:\ell_2^m\to H,\quad Ue_j=Je_j,
 \qquad C=X^*TV:H\to\ell_2^m.
\]
A Rademacher average gives
\[
 \norm U_{S_2}^2
 =\sum_{j=1}^m\norm{Je_j}^2
 =\mathbb E\norm{J(\varepsilon_1,\ldots,\varepsilon_m)}^2
 \le\norm J^2.                                                         \tag{6}
\]
Moreover, $C^*=V^*T^*X$.  By (4), the ideal property of the $2$-summing norm,
and the equality between $2$-summing and Hilbert--Schmidt norms for operators
between Hilbert spaces,
\[
 \norm C_{S_2}=\norm{C^*}_{S_2}
 \le K_G\norm V\norm T\norm X.                                        \tag{7}
\]
Thus $CU$ is trace class.  Since the diagonal map from $S_1(\ell_2^m)$ to
$\ell_1^m$ is contractive and
\[
 \ip{CUe_j}{e_j}
 =\ip{TVJe_j}{Xe_j}
 =\ip{Ty_j}{x_j}=A(x_j,y_j),
\]
we obtain from (3), (5)--(7)
\begin{align*}
 \sum_{j=1}^m|A(x_j,y_j)|
 &\le\norm{CU}_{S_1}
 \le\norm C_{S_2}\norm U_{S_2}\\
 &\le K_G^2\norm A\,
 \norm{(x_j)}_{\ell_2^w(\ell_p)}
 \norm{(y_j)}_{\ell_1^w(\ell_1)}.                                    \tag{8}
\end{align*}
This proves the $(1;2,1)$ assertion.  Since the $\ell_p$ norm of a finite
scalar family is at most its $\ell_1$ norm, the desired $(p;2,1)$ assertion
follows.

\section{Uniform Rademacher Euclidean blocks}

The counterexample uses the following explicit standard consequence of the
Khintchine inequality~\cite{Haagerup}.

\begin{lemma}\label{lem:blocks}
For every $1<a<\infty$ and every positive integer $d$, there are maps
\[
 E_{a,d}:\ell_2^d\to\ell_a^{2^d},\qquad
 Q_{a,d}:\ell_a^{2^d}\to\ell_2^d
\]
such that $Q_{a,d}E_{a,d}=I_{\ell_2^d}$ and
\[
 \norm{E_{a,d}}\le B_a,\qquad \norm{Q_{a,d}}\le B_{a'},                \tag{9}
\]
where the constants depend only on $a$.
\end{lemma}

\begin{proof}
Put $N=2^d$ and index the coordinates of $\K^N$ by
$\varepsilon\in\{-1,1\}^d$.  Define
\[
 (E_{a,d}u)_\varepsilon
 =N^{-1/a}\sum_{j=1}^d u_j\varepsilon_j,
 \qquad
 (Q_{a,d}z)_j
 =N^{-1/a'}\sum_\varepsilon z_\varepsilon\varepsilon_j.               \tag{10}
\]
Orthogonality of the signs and $1/a+1/a'=1$ give
$Q_{a,d}E_{a,d}=I$ exactly.  The upper Khintchine inequality gives the first
bound in (9).  For $v\in\ell_2^d$, H\"older's inequality followed by
Khintchine gives
\[
 |\ip{Q_{a,d}z}{v}|
 \le\norm z_a
 \left(N^{-1}\sum_\varepsilon
 \left|\sum_jv_j\varepsilon_j\right|^{a'}\right)^{1/a'}
 \le B_{a'}\norm z_a\norm v_2,
\]
which proves the second bound.
\end{proof}

\section{A compact counterexample for every \texorpdfstring{$q>1$}{q>1}}

Fix $1<q\le2$, put $s=p'>2$, and define $r>0$ by
\[
 \frac1r=\frac1p-\frac12.                                              \tag{11}
\]
Let $d_n=2^n$ and $N_n=2^{d_n}$.  Apply Lemma~\ref{lem:blocks} in each block:
\[
 E_n=E_{q,d_n},\quad Q_n=Q_{q,d_n},\qquad
 F_n=E_{s,d_n},\quad B_n=Q_{s,d_n}.                                    \tag{12}
\]
Thus $Q_nE_n=B_nF_n=I_{\ell_2^{d_n}}$, and all four operator norms are
uniformly bounded (with constants depending only on $p$ and $q$).

Partition the coordinates so that
\[
 \ell_q=\left(\bigoplus_{n\ge1}\ell_q^{N_n}\right)_q,
 \qquad
 \ell_s=\left(\bigoplus_{n\ge1}\ell_s^{N_n}\right)_s.
\]
Set
\[
 \lambda_n=d_n^{-1/(2r)},\qquad
 T=\bigoplus_{n\ge1}\lambda_nF_nQ_n:\ell_q\to\ell_s.                 \tag{13}
\]
This is bounded.  Indeed, for $z=(z_n)$ and a constant $C=C(p,q)$,
\[
 \norm{Tz}_s
 \le C\sup_n\lambda_n
       \left(\sum_n\norm{z_n}_q^s\right)^{1/s}
 \le C\sup_n\lambda_n\norm z_q,                                      \tag{14}
\]
because $q<s$.  If $T^{(N)}$ retains only the first $N$ blocks, it has finite
rank and the same estimate gives
\[
 \norm{T-T^{(N)}}\le C\sup_{n>N}\lambda_n\longrightarrow0.           \tag{15}
\]
Consequently $T$ is compact.

Let $A_T(x,y)=\ip{Ty}{x}$ on $\ell_p\times\ell_q$.  In the $n$th block, for
$1\le j\le d_n$, set
\[
 y_j^{(n)}=E_ne_j\in\ell_q,
 \qquad x_j^{(n)}=B_n^*e_j\in\ell_p.                                  \tag{16}
\]
The weak norms satisfy
\[
 \norm{(x_j^{(n)})_{j=1}^{d_n}}_{\ell_2^w(\ell_p)}
 =\norm{B_n^*}=\norm{B_n}\le B_p,                                     \tag{17}
\]
and, viewing the column map as $E_n:\ell_\infty^{d_n}\to\ell_q^{N_n}$,
\[
 \norm{(y_j^{(n)})_{j=1}^{d_n}}_{\ell_1^w(\ell_q)}
 \le\norm{E_n:\ell_2^{d_n}\to\ell_q^{N_n}}
       \norm{I:\ell_\infty^{d_n}\to\ell_2^{d_n}}
 \le B_q\sqrt{d_n}.                                                    \tag{18}
\]
On the other hand, (12)--(13) give, for every $j$,
\[
 A_T(x_j^{(n)},y_j^{(n)})
 =\lambda_n\ip{F_ne_j}{B_n^*e_j}
 =\lambda_n\ip{B_nF_ne_j}{e_j}
 =\lambda_n.                                                           \tag{19}
\]
If $A_T$ were absolutely $(p;2,1)$-summing, (17)--(19) would force
\[
 \pi_{(p;2,1)}(A_T)
 \ge\frac{\lambda_nd_n^{1/p}}{B_pB_q\sqrt{d_n}}
 =\frac1{B_pB_q}d_n^{1/(2r)}\longrightarrow\infty,                    \tag{20}
\]
where (11) and the choice of $\lambda_n$ were used.  This contradiction
proves the negative assertion and completes Theorem~\ref{thm:main}.

\section{Verification and novelty limitations}

The finite matrix identity in Lemma~\ref{lem:blocks} and the exponent in
(20) were independently checked by the included script
\texttt{code/verify\_rademacher\_blocks.py}.  The computation is not used as
proof.  The delicate mathematical points for review are the normed form (3)
of Grothendieck factorization and the application of little Grothendieck to
$V^*$ in (7).

Exact-title, equation, and exponent searches recovered the source paper and
general summability literature but no later explicit resolution of (2).  This
was a bounded search, not an exhaustive bibliographic proof.  The
classification, especially the Euclidean-block obstruction, should be checked
against operator-ideal literature using other terminology before any public
novelty claim.

\begin{thebibliography}{9}
\bibitem{BBPR}
O. Blasco, G. Botelho, D. Pellegrino, and P. Rueda,
\emph{Absolutely summing multilinear operators on $\ell_p$ spaces},
Collect. Math. 67 (2016), 519--532; arXiv:1404.1322.

\bibitem{Pisier}
G. Pisier,
\emph{Grothendieck's theorem, past and present},
Bull. Amer. Math. Soc. 49 (2012), 237--323; arXiv:1101.4195.

\bibitem{Haagerup}
U. Haagerup,
\emph{The best constants in the Khintchine inequality},
Studia Math. 70 (1981), 231--283.
\end{thebibliography}

\end{document}
